\paragraph{\trBrane{}的\trFloerCohomology。}\label{sub:floer-cohomology-for-branes}
我们需要一条初等的记号。取定一个\trCoefficients{}域 $\KK$。若 $\delta$ 是任意一个一维实\trVectorSpace，我们便把由 $\delta$ 的两个\trOrientation{}所生成的一维 $\KK$-\trVectorSpace{}记作 $|\delta|_\KK$，其\trGenerator{}之间约定：两\trGenerator{}之和为零。显然，$\RR$-\trLinear\trIsomorphism{}$c\colon\delta_1\to\delta_2$ 诱导出 $\KK$-\trLinear\trIsomorphism{}$|c|_\KK\colon|\delta_1|_\KK\to|\delta_2|_\KK$。这分别称为 $\delta$ 与 $c$ 的 $\KK$-\textbf{\trNormalization}（normalization）。\trNormalization{}对\trFamily{}同样适用：它把实\trLineBundle{}变成\trFiber{}为 $\KK$ 的\trLocalCoefficient{}系统，对它们之间的\trIsomorphism{}亦然。不妨把它看作：沿\trHomomorphism{}$\RR^{\times}\to\{\pm1\}\subset\KK^{\times}$、$t\mapsto t/|t|$ 转到相伴丛上。
升级版的\trFloerCohomology{}对 $M$ 中任何一对\trExactLagrangianBrane{}$(L_0^{\#},L_1^{\#})$ 都关联上一个有限维\trGraded{}的 $\KK$-\trVectorSpace{}$HF^*(L_0^{\#},L_1^{\#})$（取 $\mathrm{char}(\KK)=2$ 并把所有\trGraded{}部分加起来，便可由此恢复 $HF^{\mathrm{pr}}(L_0,L_1)$）。现在我们要重温 §8 中讨论过的几条基本性质，并说明如何把它们调整到有\trBraneStructure{}在场的情形。
\begin{itemize}
\item \emph{\trTriangleProduct}。存在一个\trGraded{}的、结合的、\trUnital{}的乘积
\begin{equation}\label{eq:brane-triangle-product}
HF^*(L_1^{\#},L_2^{\#})\otimes HF^*(L_0^{\#},L_1^{\#})
\longrightarrow HF^*(L_0^{\#},L_2^{\#}).
\end{equation}
\item \emph{\trBraneStructure{}的更换}。若对所涉及的每一个\trBrane{}都施行同一运算，则\trFloerCohomology{}群保持不变，\eqref{eq:brane-triangle-product} 亦然。这里的“运算”可以是：用一个 $\ZZ$ 中的常数改变\trGrading，或者用一条公共\trLineBundle{}$\beta\to M$ 的限制来张量诸\trPinStructure。略微不那么显然的是，还存在一些\trCanonical{}的\trIsomorphism（它们印证了“\trSuspension”之名）
\begin{equation}\label{eq:brane-shift-iso}
HF^*(L_0^{\#},S^{\sigma}L_1^{\#})\cong HF^{*+\sigma}(L_0^{\#},L_1^{\#}).
\end{equation}
这些\trIsomorphism{}仍与乘积\trCompatible，其含义是下图交换：
\begin{equation}\label{eq:brane-shift-product-diagram}
\begin{tikzcd}[column sep=6em,row sep=3em]
HF^*(L_1^{\#},S^{\sigma}L_2^{\#})\otimes HF^*(L_0^{\#},L_1^{\#})
\ar[r]\ar[d,"\cong"]
&
HF^*(L_0^{\#},S^{\sigma}L_2^{\#})
\ar[d,"\cong"]
\\
HF^{*+\sigma}(L_1^{\#},L_2^{\#})\otimes HF^*(L_0^{\#},L_1^{\#})
\ar[r]
&
HF^{*+\sigma}(L_0^{\#},L_2^{\#}).
\end{tikzcd}
\end{equation}
\item \emph{PSS~\trIsomorphism}。存在一个\trCanonical{}的\trIsomorphism{}$HF^*(L^{\#},L^{\#})\cong H^*(L;\KK)$，它把\eqref{eq:brane-triangle-product}（取 $L_0^{\#}=L_1^{\#}=L_2^{\#}=L^{\#}$）变成寻常的\trCupProduct。更一般地，若用任意两条\trLineBundle{}$\beta_0$、$\beta_1$ 来\trTwisted{}$L^{\#}$，则所得是 $L$ 取值于相应\trLocalCoefficient{}系统的\trCohomology：
\begin{equation}\label{eq:pss-twisted-cohomology}
HF^*(L^{\#}\otimes\beta_0,\,L^{\#}\otimes\beta_1)\cong H^*\bigl(L;\,|\beta_0^{\vee}\otimes\beta_1|_\KK\bigr).
\end{equation}
\item \emph{对偶}。对任何 $L^{\#}$，\trPoincareDuality{}与 PSS~\trIsomorphism{}给出一个映射
\begin{equation}\label{eq:brane-duality-map}
\int_{L^{\#}}\colon
HF^n\bigl(L^{\#},L^{\#}\otimes\lambda^{\mathrm{top}}(TL)\bigr)
\cong H^n\bigl(L;\,|\lambda^{\mathrm{top}}TL|_\KK\bigr)
\cong H_0(L;\KK)\to\KK.
\end{equation}
由它构造的、\trDegree{}为 $-n$ 的\trPairing
\begin{equation*}
HF^*(L_1^{\#},L_0^{\#}\otimes\lambda^{\mathrm{top}}(TL_0))\otimes HF^*(L_0^{\#},L_1^{\#})\to\KK,
\qquad
\langle a_2,a_1\rangle=\int_{L_0^{\#}}a_2a_1
\end{equation*}
\begingroup\tolerance=9999\emergencystretch=3em\hbadness=10000
是\trNondegenerate{}而且\trGraded{}对称的。\penalty0后一条性质需要稍作解释。\penalty0利用取定的\trIsomorphism{}$\lambda^{\mathrm{top}}(TL_k)\cong\xi_M|_{L_k}$，可以把 $a_2$ 看成 $HF^*(L_1^{\#},L_0^{\#}\otimes\xi_M|_{L_0})$ 中的元素，而把 $a_1$ 看成 $HF^*(L_1^{\#}\otimes\xi_M|_{L_0},L_1^{\#}\otimes\lambda^{\mathrm{top}}(TL_1))$ 中的元素。记住这一点，便有 $a_1a_2\in HF^*(L_1^{\#},L_1^{\#}\otimes\lambda^{\mathrm{top}}(TL_1))$，从而下列方程是有意义的：\par\endgroup
\begin{equation*}
\langle a_1,a_2\rangle=(-1)^{|a_1||a_2|}\langle a_2,a_1\rangle.
\end{equation*}
\item \emph{\trIsotopy\trInvariance}。从 $L_0^{\#}$ 到 $L_1^{\#}$ 的一个\trExactLagrangianBrane{}\trIsotopy{}给出 $HF^0(L_0^{\#},L_1^{\#})$ 中的一个取定元素。这些元素的性质一如\cref{sub:moving-boundary-conditions}所述；特别地，可以用它们证明 $HF^*$ 在\trIsotopy{}下不变。
\end{itemize}
我们按如下方式定义\trDonaldson--\trFukayaCategory{}$H\scr{F}(M)$：取\trExactLagrangianBrane{}为\trObject、\trFloerCohomology{}群为\trMorphism，而\trComposition{}则由\eqref{eq:brane-triangle-product} 给出。这是一个 $\KK$-\trLinear、$\ZZ$-\trGraded{}的\trCategory（若取 $\mathrm{char}(\KK)=2$ 并遗忘\trGrading，它就\trDegenerate{}为 $H\scr{F}(M)^{\mathrm{pr}}$）。由于\trPoincareDuality{}\trPairing{}的\trTwisted{}形式，它一般不是\trFrobenius\trCategory，除非 $\xi_M$ 是平凡的。
